% GENERATED FILE. Edit canonical lessons, then run npm run generate:books.
% canonical-source-hash: fnv1a64:16a78987f6037ea2
% canonical-lessons: KA-C406-visvavidyalaya, KA-C406-padavi, KA-C406-upanyasa, KA-C406-pradhyapaka, KA-C406-bodhisu

\chapter{At the University}
\label{ch:ka-at-the-university}
\hlchaptermodality{406}

\begin{chapteropening}
\textbf{By the end of this chapter:} \emph{I can say a university, a degree, a lecture, a professor and to instruct.}

Complete the last lesson of chapter 406: I can say a university, a degree, a lecture, a professor and to instruct.
\end{chapteropening}

\section[\texorpdfstring{\kn{ವಿಶ್ವವಿದ್ಯಾಲಯ} (viśvavidyālaya)}{ವಿಶ್ವವಿದ್ಯಾಲಯ (viśvavidyālaya)}]{\kn{ವಿಶ್ವವಿದ್ಯಾಲಯ} (viśvavidyālaya) — a university}

\label{lesson:KA-C406-visvavidyalaya}



\begin{quote}
Before the new one: say the Kannada for a compromise, then the Kannada for to incite.
\end{quote}

\subsection*{You'll want to know: \kn{ವಿಶ್ವವಿದ್ಯಾಲಯ}}
\textbf{\kn{ವಿಶ್ವವಿದ್ಯಾಲಯ}} — \emph{viśvavidyālaya} — \textquotedblleft{}a university\textquotedblright{}.

\begin{quote}
\kn{ನನ್ನ} \kn{ಅಕ್ಕ} \kn{ವಿಶ್ವವಿದ್ಯಾಲಯದಲ್ಲಿ} \kn{ಓದುತ್ತಾಳೆ}.
\end{quote}

\begin{quote}
\emph{nanna akka viśvavidyālayadalli ōduttāḷe.}
\end{quote}

\begin{quote}
My elder sister studies at the university.
\end{quote}

Plural: \kn{ವಿಶ್ವವಿದ್ಯಾಲಯಗಳು}, \emph{viśvavidyālayagaḷu}.

\kn{ವಿಶ್ವ} (world) + \kn{ವಿದ್ಯಾ} (learning) + \kn{ಆಲಯ} (house).

\begin{cousinweb}
\textbf{\kn{ವಿಶ್ವವಿದ್ಯಾಲಯ}} (\emph{viśvavidyālaya}) is a modern word built from Sanskrit \emph{viśva}, \textquotedblleft{}all, the world\textquotedblright{} (CDIAL 11962), \emph{vidyā}, \textquotedblleft{}learning\textquotedblright{} (CDIAL 11741), and \emph{ālaya}, \textquotedblleft{}a house\textquotedblright{} (CDIAL 1366).
\end{cousinweb}

\begin{grammarlens}[title={What you've built}]
One more word for giving reasons and opinions.
\end{grammarlens}

\subsection*{Your turn}
\begin{itemize}
\raggedright
  \item \emph{Say it:} \emph{viśvavidyālaya}
  \item \emph{Say it:} \emph{viśvavidyālaya}, once more
  \item \emph{Say it:} \emph{nanna akka viśvavidyālayadalli ōduttāḷe.} — My elder sister studies at the university.
  \item \emph{Recall:} say the Kannada for a compromise, then the Kannada for to incite, then say \emph{viśvavidyālaya} again
\end{itemize}

\begin{tcolorbox}[breakable,colback=teal!4,colframe=teal!35!black,title={Before you move on}]
What does \kn{ವಿಶ್ವವಿದ್ಯಾಲಯ} mean? (\textquotedblleft{}A university\textquotedblright{}.) Say it once more.
\end{tcolorbox}
\section[\texorpdfstring{\kn{ಪದವಿ} (padavi)}{ಪದವಿ (padavi)}]{\kn{ಪದವಿ} (padavi) — a degree; a rank}

\label{lesson:KA-C406-padavi}



\begin{quote}
Before the new one: say the Kannada for to incite, then the Kannada for a university.
\end{quote}

\subsection*{You'll want to know: \kn{ಪದವಿ}}
\textbf{\kn{ಪದವಿ}} — \emph{padavi} — \textquotedblleft{}a degree; a rank\textquotedblright{}.

\begin{quote}
\kn{ಅವನು} \kn{ಕಳೆದ} \kn{ವರ್ಷ} \kn{ಪದವಿ} \kn{ಮುಗಿಸಿದನು}.
\end{quote}

\begin{quote}
\emph{avanu kaḷeda varṣa padavi mugisidanu.}
\end{quote}

\begin{quote}
He finished his degree last year.
\end{quote}

Also a post or rank: \kn{ಉನ್ನತ} \kn{ಪದವಿ}, a high post.

\begin{cousinweb}
\textbf{\kn{ಪದವಿ}} (\emph{padavi}) is borrowed from Sanskrit \emph{padavī}, \textquotedblleft{}a path; a station, a position\textquotedblright{} (Monier-Williams).
\end{cousinweb}

\begin{grammarlens}[title={What you've built}]
One more word for giving reasons and opinions.
\end{grammarlens}

\subsection*{Your turn}
\begin{itemize}
\raggedright
  \item \emph{Say it:} \emph{padavi}
  \item \emph{Say it:} \emph{padavi}, once more
  \item \emph{Say it:} \emph{avanu kaḷeda varṣa padavi mugisidanu.} — He finished his degree last year.
  \item \emph{Recall:} say the Kannada for to incite, then the Kannada for a university, then say \emph{padavi} again
\end{itemize}

\begin{tcolorbox}[breakable,colback=teal!4,colframe=teal!35!black,title={Before you move on}]
What does \kn{ಪದವಿ} mean? (\textquotedblleft{}A degree; a rank\textquotedblright{}.) Say it once more.
\end{tcolorbox}
\section[\texorpdfstring{\kn{ಉಪನ್ಯಾಸ} (upanyāsa)}{ಉಪನ್ಯಾಸ (upanyāsa)}]{\kn{ಉಪನ್ಯಾಸ} (upanyāsa) — a lecture}

\label{lesson:KA-C406-upanyasa}



\begin{quote}
Before the new one: say the Kannada for a university, then the Kannada for a degree.
\end{quote}

\subsection*{You'll want to know: \kn{ಉಪನ್ಯಾಸ}}
\textbf{\kn{ಉಪನ್ಯಾಸ}} — \emph{upanyāsa} — \textquotedblleft{}a lecture\textquotedblright{}.

\begin{quote}
\kn{ಇಂದಿನ} \kn{ಉಪನ್ಯಾಸ} \kn{ಒಂದು} \kn{ಗಂಟೆ} \kn{ನಡೆಯಿತು}.
\end{quote}

\begin{quote}
\emph{indina upanyāsa ondu gaṇṭe naḍeyitu.}
\end{quote}

\begin{quote}
Today's lecture went on for an hour.
\end{quote}

\kn{ಉಪನ್ಯಾಸಕ}, a lecturer.

\begin{cousinweb}
\textbf{\kn{ಉಪನ್ಯಾಸ}} (\emph{upanyāsa}) is borrowed from Sanskrit \emph{upanyāsa}, \textquotedblleft{}putting down near; a statement\textquotedblright{} (Monier-Williams).
\end{cousinweb}

\begin{grammarlens}[title={What you've built}]
One more word for giving reasons and opinions.
\end{grammarlens}

\subsection*{Your turn}
\begin{itemize}
\raggedright
  \item \emph{Say it:} \emph{upanyāsa}
  \item \emph{Say it:} \emph{upanyāsa}, once more
  \item \emph{Say it:} \emph{indina upanyāsa ondu gaṇṭe naḍeyitu.} — Today's lecture went on for an hour.
  \item \emph{Recall:} say the Kannada for a university, then the Kannada for a degree, then say \emph{upanyāsa} again
\end{itemize}

\begin{tcolorbox}[breakable,colback=teal!4,colframe=teal!35!black,title={Before you move on}]
What does \kn{ಉಪನ್ಯಾಸ} mean? (\textquotedblleft{}A lecture\textquotedblright{}.) Say it once more.
\end{tcolorbox}
\section[\texorpdfstring{\kn{ಪ್ರಾಧ್ಯಾಪಕ} (prādhyāpaka)}{ಪ್ರಾಧ್ಯಾಪಕ (prādhyāpaka)}]{\kn{ಪ್ರಾಧ್ಯಾಪಕ} (prādhyāpaka) — a professor}

\label{lesson:KA-C406-pradhyapaka}



\begin{quote}
Before the new one: say the Kannada for a degree, then the Kannada for a lecture.
\end{quote}

\subsection*{You'll want to know: \kn{ಪ್ರಾಧ್ಯಾಪಕ}}
\textbf{\kn{ಪ್ರಾಧ್ಯಾಪಕ}} — \emph{prādhyāpaka} — \textquotedblleft{}a professor\textquotedblright{}.

\begin{quote}
\kn{ಆ} \kn{ಪ್ರಾಧ್ಯಾಪಕರು} \kn{ಗಣಿತ} \kn{ಕಲಿಸುತ್ತಾರೆ}.
\end{quote}

\begin{quote}
\emph{ā prādhyāpakaru gaṇita kalisuttāre.}
\end{quote}

\begin{quote}
That professor teaches mathematics.
\end{quote}

Plural: \kn{ಪ್ರಾಧ್ಯಾಪಕರು}, \emph{prādhyāpakaru}.

Respectful plural verb: \kn{ಪ್ರಾಧ್ಯಾಪಕರು} \kn{ಬಂದರು}.

\begin{cousinweb}
\textbf{\kn{ಪ್ರಾಧ್ಯಾಪಕ}} (\emph{prādhyāpaka}) is a modern word built on Sanskrit \emph{adhyāpaka}, \textquotedblleft{}a teacher\textquotedblright{} (CDIAL 278), with \emph{pra-}, \textquotedblleft{}forward, chief\textquotedblright{}.
\end{cousinweb}

\begin{grammarlens}[title={What you've built}]
One more word for giving reasons and opinions.
\end{grammarlens}

\subsection*{Your turn}
\begin{itemize}
\raggedright
  \item \emph{Say it:} \emph{prādhyāpaka}
  \item \emph{Say it:} \emph{prādhyāpaka}, once more
  \item \emph{Say it:} \emph{ā prādhyāpakaru gaṇita kalisuttāre.} — That professor teaches mathematics.
  \item \emph{Recall:} say the Kannada for a degree, then the Kannada for a lecture, then say \emph{prādhyāpaka} again
\end{itemize}

\begin{tcolorbox}[breakable,colback=teal!4,colframe=teal!35!black,title={Before you move on}]
What does \kn{ಪ್ರಾಧ್ಯಾಪಕ} mean? (\textquotedblleft{}A professor\textquotedblright{}.) Say it once more.
\end{tcolorbox}
\section[\texorpdfstring{\kn{ಬೋಧಿಸು} (bōdhisu)}{ಬೋಧಿಸು (bōdhisu)}]{\kn{ಬೋಧಿಸು} (bōdhisu) — to instruct; to preach}

\label{lesson:KA-C406-bodhisu}



\begin{quote}
Before the new one: say the Kannada for a lecture, then the Kannada for a professor.
\end{quote}

\subsection*{You'll want to know: \kn{ಬೋಧಿಸು}}
\textbf{\kn{ಬೋಧಿಸು}} — \emph{bōdhisu} — \textquotedblleft{}to instruct; to preach\textquotedblright{}.

\begin{quote}
\kn{ಅವರು} \kn{ಮೂವತ್ತು} \kn{ವರ್ಷ} \kn{ವಿಜ್ಞಾನ} \kn{ಬೋಧಿಸಿದರು}.
\end{quote}

\begin{quote}
\emph{avaru mūvattu varṣa vijñāna bōdhisidaru.}
\end{quote}

\begin{quote}
They taught science for thirty years.
\end{quote}

Past (I …): \kn{ಬೋಧಿಸಿದೆ}, \emph{bōdhiside}.

More formal than \kn{ಕಲಿಸು}. \kn{ಬೋಧಕ}, a teacher.

\begin{cousinweb}
\textbf{\kn{ಬೋಧಿಸು}} (\emph{bōdhisu}) is built on Sanskrit \emph{bodhana}, \textquotedblleft{}arousing, awakening\textquotedblright{} (Monier-Williams), from \emph{budh}, \textquotedblleft{}to wake, to know\textquotedblright{}, with the Kannada verb ending -\emph{isu}.
\end{cousinweb}

\begin{grammarlens}[title={What you've built}]
One more word for giving reasons and opinions.
\end{grammarlens}

\subsection*{Your turn}
\begin{itemize}
\raggedright
  \item \emph{Say it:} \emph{bōdhisu}
  \item \emph{Say it:} \emph{bōdhisu}, once more
  \item \emph{Say it:} \emph{avaru mūvattu varṣa vijñāna bōdhisidaru.} — They taught science for thirty years.
  \item \emph{Recall:} say the Kannada for a lecture, then the Kannada for a professor, then say \emph{bōdhisu} again
\end{itemize}

\begin{tcolorbox}[breakable,colback=teal!4,colframe=teal!35!black,title={Before you move on}]
What does \kn{ಬೋಧಿಸು} mean? (\textquotedblleft{}To instruct; to preach\textquotedblright{}.) Say it once more.
\end{tcolorbox}
